Según la prensa consultada, el 8 de octubre de 2026 un flujo de barro descendió desde la quebrada de San Carlos de Apoquindo (Las Condes, Santiago) sobre una ladera quemada nueve meses antes. Se ensayó su reconstrucción con un modelo bidimensional de aguas someras con reología de barro, sobre un relieve de \SI{30}{m} remuestreado a una retícula de \SI{4}{m} con 42.018 edificios, usando solo datos públicos y un registro probatorio de cada fuente consultada. Se completaron 84 simulaciones: 19 de cribado, 41 de ciudad completa, 5 de cuadratura y 19 microsimulaciones del campus. Con unos \SI{144}{mm} acumulados según lo informado por la alcaldesa y con la isoterma de \SI{0}{\celsius} informada por el CR2 por sobre la cumbre de la cuenca (área calculada: \SI{10,1}{km^2}), una estimación independiente que adopta la saturación del suelo informada sitúa el volumen de mezcla en unos \SI{295000}{m^3}, del mismo orden de magnitud que el requerido por el alcance simulado. La hipótesis seleccionada para la prueba reservada (z02) ---cauce principal, brazos sur y lluvia sobre la ladera--- reproduce los corredores de Camino El Alba, General Blanche y Quebrada Honda que delimitan la zona afectada oficial, el alcance hasta Padre Hurtado y cerca de \SI{1}{m} de barro en la Clínica UC Christus, con un índice exploratorio de compatibilidad (IEC) de 72/100. Este índice compara hipótesis y no mide precisión predictiva. Una prueba con observaciones reservadas, definida antes de examinar los resultados, resultó desfavorable: el flujo simulado alcanza 19 de 49 equipamientos afectados oficiales (\SI{39}{\percent}, frente al \SI{70}{\percent} fijado), porque el modelo deja sin barro partes de la zona afectada oficial. En los ensayos posteriores de ajuste, la entrada de barro en las casas representadas y el supuesto de agua turbia en el campus y La Plaza Sur, junto con barro en la rama principal, elevan el IEC verificado con paso fino a 76--82/100. La hipótesis integrada obtiene 79--82/100: 82 con la lluvia supuesta de \SI{35}{mm/h} y paso fino (z32), 81 con la lluvia medida en Río Molina con su forma de pulso (p34b, corregida) y 79 con esa lluvia repartida en \SI{4}{h} (v33). La severidad del incendio, estimada con Sentinel-2, muestra señal de quema en el \SI{82}{\percent} de las cuencas, en su mayoría con severidad baja o moderada-baja. Con los coeficientes supuestos para el escenario sin incendio, el balance calcula alrededor del \SI{62}{\percent} del volumen de barro y el \SI{43}{\percent} de los sólidos del escenario con incendio. En el modelo, un obstáculo supuesto de troncos al final del cauce principal no desvía el flujo hacia el sur. En escenarios hipotéticos, la inundación de Quebrada Honda se reproduce si el sumidero que representa el colector de Francisco Bulnes Correa no retira todo el flujo calculado en su embocadura, de unos \SI{65}{m^3/s}. Una capacidad supuesta de \SI{30}{m^3/s} deja flujo en superficie. Estos valores no describen la obra real. La retención en las casas representadas mejora la coincidencia con la zona afectada oficial, pero ninguna configuración reproduce a la vez esa extensión y el metro informado en la Clínica UC Christus. Las microsimulaciones y la simulación de ciudad con malla fina en el campus (m36) reproducen la entrada por el claro al patio hundido; m36 calcula \SI{0,17}{m} en el patio y moja el \SI{58}{\percent} de La Plaza Sur. El estudio se cerró el 13 de octubre de 2026; sus resultados de ajuste no tienen validación independiente.
